\section{One higher-rank component in each row}\label{sec:higher-rank}
The rank-one condition can be relaxed at one port in every row, without
bounding that port's rank or any original coordinate width. The graph
hypotheses and intrinsic capacity convention remain those of
\cref{sec:rank-one}.

\begin{theorem}[H: one higher-rank port]\label{thm:one-high}
Let $C\leq\prod_v A_v$ have prescribed linear recoveries
$x_v=\sum_u T_{vu}x_u$, with $T_{vu}=0$ off a finite simple triangle-free
subcubic graph. Suppose, for each $v$, at most one map $T_{vu}$ has rank
greater than one. Then
\[
 2\dim C\leq\sum_v\dim\pi_v(C)\leq\sum_v\dim A_v.
\]
For the corresponding block operator $B$ of total scalar size $N$, this
gives $2\rank(I-B)\geq N$.
\formal{HigherRankCodes.lean}{LowRank.one_high_code_half,LowRank.one_high_code_capacity,LowRank.one_high_block_rank}
\end{theorem}
\begin{outline}
Use the coordinate-dual arrangement and induction on ambient dimension,
as in \cref{thm:rank-one}. Local sections preserve the one-higher-rank
condition in every quotient. A counterexample would again have
$f(S)\geq1$ for every nonzero test space, and distinct coordinate spaces
would be disjoint.

Each coordinate is a line or plane. For a row with a higher-rank component
$R_{vu}$, the remainder $x\mapsto x-R_{vu}x$ is injective: a vector in its
kernel would lie in both $U_v$ and $U_u$. Its image is spanned by the at
most two remaining rank-one ports, so $\dim U_v\leq2$. The previous
three-line argument handles rows without a higher-rank port. Thus every
higher-rank component in this hypothetical core is a rank-two
isomorphism between planes. Call its directed edge a high arc.

A high arc $v\to u$ forces a high successor $u\to w$ with $w\ne v$.
Otherwise $U_v+U_u$ has dimension four and contains the two endpoint
planes and four exterior input lines at distinct vertices. Their eight
coordinate dimensions contradict positive deficit. This argument
includes a low-rank or zero reverse port; reciprocity is not assumed.

Along a high path $p\to a\to b\to c$, let $z$ and $y$ be the third
graph neighbors of $a$ and $b$. The space
\[
 S=U_a+U_b+\operatorname{im}R_{by}
\]
has dimension at most five and contains the planes $U_a,U_b,U_c$,
as well as input lines at $p,z,y$. Triangle-freeness makes these six
vertices distinct. Their total contribution is nine, forcing $\dim S=5$,
$f(S)=1$, and exactly those coordinate intersections.

The next successor plane lies in $S$ plus the next third-neighbor line.
That line is outside $S$. Adding it grows the space by one dimension
and increases its own coordinate intersection. Positive deficit forces
the successor plane to have already met $S$ in a line. The exact support
just proved then forces either a high four-cycle or an edge from $a$
to that successor.

A high four-cycle is impossible. Starting from its first tight triple
space, adjoining the next third-neighbor line grows dimension by one.
It gains one dimension at the remaining cycle plane and at each of two
distinct external coordinates. If an external vertex coincides with
one from the previous triple, the neighboring triple's exact intersection
shows that its new line is independent of the old line. Thus there are
three coordinate gains for two units of dimension budget, contradicting
positive deficit.

Seven consecutive high arcs now suffice for the final contradiction.
The forced distance-three edges and maximum degree three make the last
vertex equal the second. The six intervening distinct vertices form
$K_{3,3}$, with every available graph edge used. Recovery makes all six
plane coordinates lie in the span of one three-vertex part. That span
has dimension at most six but at least twelve coordinate-intersection
dimensions, again contradicting positive deficit. Every high arc could
be extended to such a path, so there are none. The rank-one arrangement
bound finishes the induction. The code and block-rank conclusions use
the same dualization and rank--nullity as before.
\end{outline}

The arbitrary-width arrangement theorem is
\formal{HigherRankBound.lean}{LowRank.one_high_arrangement_half}
Its two final exclusions are
\formal{HigherRankFourCycle.lean}{LowRank.HighCore.no_high_four_cycle}
\formal{HigherRankExclusion.lean}{LowRank.HighCore.no_seven_path,LowRank.HighCore.no_high}
The proof does not require a classification of the full high-arc graph.
In particular, it does not assume that high arcs have distinct predecessors.

\begin{corollary}[Exceptional rows for Theorem H]\label{cor:one-high-exceptions}
In the arrangement version of \cref{thm:one-high}, allow rows at
$D\subseteq V$ to have arbitrarily many higher-rank components. With
$H=\sum_{v\in D}U_v$ and $\Delta=2\dim W-\sum_v\dim U_v$, the exact
payment remains
\[
 \Delta\leq2\dim H-\sum_v\dim(U_v\cap H)
 \leq\dim H\leq\sum_{v\in D}\dim U_v.
\]
\formal{HigherRankCodes.lean}{LowRank.one_high_exception_defect}
\formal{RankOneCodes.lean}{LowRank.exceptional_span_defect}
\end{corollary}
\begin{outline}
Quotient by $H$. Exceptional coordinates become zero; all other rows retain
at most one higher-rank component. Apply the arrangement theorem and
the same exact intersection accounting as in
\cref{thm:rank-one-exceptions}.
\end{outline}

For a construction exceeding half capacity, some rows must have at least
two higher-rank ports, and the span of those exceptional coordinates must
pay the excess. This does not imply a matching-number bound for arbitrary
vector codes, nor a small exceptional span for native computations.
The reversible-rule and one-higher-rank hypotheses are separate sufficient
conditions; neither is asserted for unrestricted linear or nonlinear
recovery.
